%% kapa: exam=1 round=36 year=2025 period=1 form=A subject=economics subject_ko=경제학원론 seq=14 number=54
%% source: Q-Net 감정평가사 기출문제 (cst00309) artlSeq=5231525 file=raw/1008/5231525/2025년 제36회 감정평가사 1차시험 1교시 문제지 원본.pdf
%% conversion: kordoc 4.24.51 -> tools/gen1.py (auto); review=proofread-sonnet; flags=none
\begin{question}{54}
생산요소시장에서 경제적 지대(economic rent)에 관한 설명으로 옳지 않은 것은?
\begin{choices}
  \item 생산요소의 기회비용을 초과해 추가로 지불되는 보수로 해석할 수 있다.
  \item 전용수입(transfer earnings)은 생산요소를 현재 수준으로 유지하기 위한 것과 관련된 기회비용이다.
  \item 공급이 비탄력적일수록 경제적 지대는 커진다.
  \item 희소성이 큰 생산요소일수록 경제적 지대가 크다.
  \item 공급이 완전탄력적일 경우 전용수입이 발생하지 않는다.
\end{choices}
\end{question}
\answer{5}
